% dh-tracking-listing.tex -- standalone listing for the Part III arXiv submission.
%
% Two contrasting kernels for doubly heterogeneous media, taken verbatim (with
% marked elisions) from outram-mc-libs:
%   Listing 1  crates/outram-mc-libs/src/pebble_beds/delta_tracking.rs
%   Listing 2  crates/outram-mc-libs/src/stochastic/{cls,scls}.rs
%
% Compiles on its own:  pdflatex dh-tracking-listing.tex
% To drop into a paper instead, copy the \lstdefinelanguage block and the two
% lstlisting environments; nothing else here is load-bearing.
%
% GPL-3.0-only, (C) 2026 Theodore Ong and the outram-park contributors.

\documentclass[11pt]{article}
\usepackage[margin=1in]{geometry}
\usepackage{listings}
\usepackage{xcolor}
\usepackage[T1]{fontenc}

\definecolor{kw}{RGB}{0,0,160}
\definecolor{cm}{RGB}{95,125,95}
\definecolor{st}{RGB}{150,60,20}

\lstdefinelanguage{Rust}{
  morekeywords={as,break,const,continue,crate,else,enum,false,fn,for,if,impl,
    in,let,loop,match,mod,move,mut,pub,ref,return,self,Self,struct,super,trait,
    true,type,unsafe,use,where,while,f64,u32,u64,usize,bool,Option,Some,None},
  sensitive=true,
  morecomment=[l]{//},
  morecomment=[s]{/*}{*/},
  morestring=[b]",
}

\lstset{
  language=Rust,
  basicstyle=\ttfamily\footnotesize,
  keywordstyle=\color{kw}\bfseries,
  commentstyle=\color{cm}\itshape,
  stringstyle=\color{st},
  numbers=left,
  numberstyle=\tiny\color{gray},
  numbersep=8pt,
  frame=single,
  rulecolor=\color{gray},
  breaklines=true,
  breakatwhitespace=true,
  showstringspaces=false,
  columns=fullflexible,
  captionpos=b,
  xleftmargin=2em,
  framexleftmargin=1.5em,
}

\begin{document}

\begin{lstlisting}[caption={Delta (Woodcock) tracking. The flight is sampled on a
  majorant $\Sigma_{\mathrm{maj}} \geq \Sigma_t(\mathbf{r})$ everywhere, so the
  neutron crosses material boundaries without stopping at them; the resulting
  bias is removed exactly by rejecting a fraction $1 - \Sigma_t/\Sigma_{\mathrm{maj}}$
  of collisions as virtual. No geometric approximation is made.},
  label={lst:delta}]
/// Flight distance [cm] on the majorant: s = -ln(xi) / Sigma_maj.
pub fn sample_delta_distance(majorant: f64, seed: &mut u64) -> f64 {
    if majorant <= 0.0 { return f64::INFINITY; }
    -prn(seed).max(f64::MIN_POSITIVE).r_ln() / majorant
}

/// Accept a real collision with probability Sigma_t(local) / Sigma_maj.
pub fn classify_collision(sigma_t: f64, majorant: f64, seed: &mut u64) -> DeltaEvent {
    if majorant <= 0.0 { return DeltaEvent::Virtual; }
    let p_real = (sigma_t / majorant).clamp(0.0, 1.0);
    if prn(seed) < p_real { DeltaEvent::Real } else { DeltaEvent::Virtual }
}

// ... drive the neutron to its next REAL collision:
loop {
    let s = sample_delta_distance(maj, seed);
    pos = pos + Position::new(dir.u * s, dir.v * s, dir.w * s);
    total += s;
    match sigma_t_at(pos) {                  // None => left the tracking region
        None => return DeltaFlight { position: pos, escaped: true, .. },
        Some(sigma_t) => match classify_collision(sigma_t, maj, seed) {
            DeltaEvent::Real    => return DeltaFlight { position: pos, .. },
            DeltaEvent::Virtual => { virtuals += 1; }   // no physics; fly on
        },
    }
}
\end{lstlisting}

\vspace{1em}

\begin{lstlisting}[caption={Semi-implicit chord-length sampling (SCLS). No
  geometry is stored; the matrix gap to the next inclusion is drawn from the
  Markovian chord distribution $\ell = -\langle\ell\rangle \ln \xi$ with
  $\langle\ell_{\mathrm{matrix}}\rangle = \tfrac{4r}{3}\,(1-p_f)/p_f$. The
  \emph{semi-implicit} part is lines 9--15: inclusions already met are retained
  inside a moving window, and if a remembered sphere lies nearer along the ray
  than the freshly sampled chord, the neutron re-meets that particle instead of
  an independent one. This restores exactly the correlation that plain CLS
  discards on back-scatter.},
  label={lst:scls}]
/// Chord length [cm]: inverse-transform sampling of an exponential.
pub fn sample_chord(mean_chord: f64, seed: &mut u64) -> f64 { /* -mean * ln(xi) */ }

/// Mean matrix chord for spherical inclusions at packing fraction pf.
pub fn matrix_mean_chord_length(radius: f64, pf: f64) -> f64 {
    mean_chord_length_sphere(radius) * (1.0 - pf) / pf      // (4r/3)(1-pf)/pf
}

// ... inside the flight reconstruction, when the neutron is in the matrix:
let d_sample = sample_chord(mean_matrix, seed);   // memoryless matrix gap ...
let mut d_incl = d_sample;                        // ... unless a REMEMBERED
let mut hit: Option<(Position, f64)> = None;      //     inclusion lies closer
for h in &self.histories {
    if let Some(d) = ray_sphere_entry(pos, u, h.center, h.radius, EPS) {
        if d < d_incl { d_incl = d; hit = Some((h.center, h.radius)); }
    }
}
if t + d_incl <= len {
    t += d_incl;
    pos = pos + u_pos * d_incl;
    fl.in_inclusion = true;
    fl.current_sphere = Some(match hit {
        Some(existing) => existing,                    // re-meet a known particle
        None => {                                      // or place a new one and
            let c = fresh_inclusion_center(pos, u, radius, seed);   // remember it
            self.remember_inclusion(ParticleHistory::new(c, radius, inclusion_mat));
            (c, radius)
        }
    });
}
self.advance_to(position);   // re-centre the retention window; cull what left it
\end{lstlisting}

\end{document}
